\section{Part 2：如何做深度 RL 研究}

\subsection{研究的心态与方法论}

\begin{importantbox}{Chelsea Finn 的研究建议}
\begin{enumerate}
    \item \textbf{从问题出发}：先理解要解决什么问题，再考虑用什么方法。不要"拿着锤子找钉子"
    \item \textbf{先理解再创新}：深入理解现有方法的优缺点，创新自然产生
    \item \textbf{简单优先}：从最简单的方法/实验开始，逐步增加复杂性
    \item \textbf{建立直觉}：用简单环境建立对算法行为的直觉，然后推广到复杂场景
\end{enumerate}
\end{importantbox}

\subsection{实验设计}

\begin{knowledgebox}{RL 实验的最佳实践}
\begin{itemize}
    \item \textbf{控制变量}：每次只改变一个因素
    \item \textbf{多种子运行}：至少 3--5 个随机种子
    \item \textbf{有意义的基线}：选择强基线而非稻草人
    \item \textbf{消融实验}：验证每个设计选择的贡献
    \item \textbf{诊断工具}：监控学习曲线、值函数估计误差、策略熵等
\end{itemize}
\end{knowledgebox}

\begin{figure}[H]
\centering
\includegraphics[width=0.9\textwidth]{slides-images/slide-15.jpg}
\caption{课程把“如何扩展 RL”具体化为 batch-online、价值函数精度和更新频率三个实验设计问题，这些问题决定算法能否真正扩展到长时程和大模型场景。}
\end{figure}

\subsection{调试 RL 算法}

RL 算法的调试比监督学习困难得多。讲者给出的调试策略：

\begin{enumerate}
    \item \textbf{在简单环境中验证}：先确认算法在 CartPole/Pendulum 上正确工作
    \item \textbf{检查奖励曲线}：奖励应该单调（或大致单调）增长
    \item \textbf{检查值函数}：值函数估计应该接近真实回报
    \item \textbf{可视化策略行为}：观察智能体实际在做什么
    \item \textbf{检查梯度}：确保梯度没有爆炸或消失
\end{enumerate}

\begin{warningbox}{RL 调试的常见陷阱}
\begin{itemize}
    \item 奖励 scale 不当导致值函数发散
    \item 折扣因子 $\gamma$ 选择不当
    \item 数据归一化/标准化遗漏
    \item exploration 不足导致策略收敛到局部最优
    \item bug 可能被随机性掩盖——即使有 bug，RL 有时也能学到"还行"的策略
\end{itemize}
\end{warningbox}

\subsection{如何选择研究方向}

讲者对选择研究方向提供了建议：
\begin{itemize}
    \item 关注\textbf{真正重要的问题}而非热门话题
    \item 寻找理论与实践的\textbf{差距}
    \item 跨领域思考（NLP $\times$ RL、robotics $\times$ RL）
    \item 与从业者交流了解真实需求
\end{itemize}

\subsection{本章小结}

好的 RL 研究需要清晰的问题定义、严谨的实验设计和有效的调试策略。简单优先、控制变量是核心方法论。

